% Combined midterm report. Compile from the repository root:
%   latexmk -pdf report.tex
% Course limit: 5 pages including figures. References may continue past that.

\documentclass[10pt,twocolumn,letterpaper]{article}

\usepackage[pagenumbers]{midterm/report/cvpr}
%% This file contains a number of tweaks that are typically applied to the main document.
%% They are not enabled by default, but can be enabled by uncommenting the relevant lines.

%%
%% Packages used by this report that cvpr.sty does not already load
%%
\usepackage{multirow}

%%
%% Inline annotations; for predefined colors, refer to "dvipsnames" in the xcolor package:
%% https://tinyurl.com/overleaf-colors
%%
\newcommand{\red}[1]{{\color{red}#1}}
\newcommand{\todo}[1]{{\color{red}#1}}
\newcommand{\TODO}[1]{\textbf{\color{red}[TODO: #1]}}
%%
%% disable for camera ready / submission by uncommenting these lines  
%%
% \renewcommand{\TODO}[1]{}
% \renewcommand{\todo}[1]{#1}

%%
%% work harder in optimizing text layout. Typically shrinks text by 1/6 of page, enable
%% it at the very end of the writing process, when you are just above the page limit
%%
\usepackage{microtype}

%%
%% fine-tune paragraph spacing
%%
% \renewcommand{\paragraph}[1]{\vspace{.5em}\noindent\textbf{#1.}}

%%
%% globally adjusts space between figure and caption
%%
% \setlength{\abovecaptionskip}{.5em}


%%
%% Allows "the use of \paper to refer to the project name"
%% with automatic management of space at the end of the word
%%
% \usepackage{xspace}
% \newcommand{\paper}{ProjectName\xspace}

%%
%% Commonly used math definitions
%%
% \DeclareMathOperator*{\argmin}{arg\,min}
% \DeclareMathOperator*{\argmax}{arg\,max}

%%
%% Tigthen underline
%%
% \usepackage{soul}
% \setuldepth{foobar}

\definecolor{cvprblue}{rgb}{0.21,0.49,0.74}
\usepackage[pagebackref,breaklinks,colorlinks,allcolors=cvprblue]{hyperref}

\def\paperID{*****}
\def\confName{CVPR}
\def\confYear{2026}

\title{Semi-Supervised Labels and Outcome Prediction\\
for Clinical Trials\\
\vspace{0.35em}
\large CS5785 / ORIE5750 / ECE5414 --- Midterm Project}

\author{Anirudha Hadagali\\
Cornell University
\and
Michael Mazzola\\
Cornell University
\and
Tarun Mandalapu\\
Cornell University
}

\begin{document}
\maketitle

\begin{abstract}
Part~1 asks why a trial stopped early. Only 157 of 1,171 training trials have
visible reason labels. We keep the 150 reason embeddings and discard eligibility
text, diseases, drugs, and start year, because those inputs did not raise the
validation score. A Gaussian mixture only helps decide which unlabeled trials
are safe to learn from. Logistic regression, cosine neighbors, and shrinkage
discriminant analysis then name the reasons. The reported Kaggle sample-wise F1
is 0.89636.

Part~2 asks whether a trial meets its primary endpoint. The score is log loss,
and the test trials start later than the training trials. We upweight recent
trials, average in a neighbor model that also prefers recent matches, and mix
in 30\% of an earlier text-and-tree ensemble that errs on different trials. The
reported Kaggle log loss is 0.58465. A blend that looked better on our
historical checks scored 0.58492, so we kept the first result.
\end{abstract}

\noindent\textbf{Team responsibilities.}
All three authors worked on both parts and on this report.
\textbf{Michael Mazzola} led the Part~1 model.
\textbf{Tarun Mandalapu} built the first Part~2 classifier, and
\textbf{Anirudha Hadagali} led the later Part~2 tuning.
Tarun Mandalapu drafted the longer project notes that this summary condenses.

\section{Part 1: stopping reasons}
\label{sec:p1}

\subsection{Data and the score}

The files hold 1,561 trials that stopped early~\cite{fu2022hint,clinicaltrials}.
Training has 1,171 trials: 157 with real labels and 1,014 with every reason
marked $-100$, which means unknown. Validation has 78 labeled trials. Test has
312 later trials, from 2010 to 2016. A trial can carry more than one of Safety,
Funding, Efficacy, Recruitment, and Other. Among the 235 real labels,
Recruitment appears on 64\% and Funding on 5\%. The labeled trials are mostly
early, so those rates will not match the test set.

The selection score is sample-wise F1, the average of a per-trial F1. We also
report F1 for each reason and the average of those five scores. One validation
trial moves sample-wise F1 by about 1.3 points, so we rechecked every choice
with eight repeats of five grouped folds, once on the 157 training labels and
once on all 235, plus three time-ordered windows.

\subsection{What the model does}

Each trial is represented by 150 numbers from its stopping-reason text. We
standardize them using only rows the current split may learn from. IDs are used
to line up rows and are not features. Adding start year, eligibility text,
disease names, pathways, or drug fingerprints~\cite{rogers2010extended,rdkit,milacic2024reactome}
never raised validation sample-wise F1 on the earlier self-training model
(0.893 for embeddings alone, 0.891 with year, 0.876 with eligibility or disease
text, 0.889 with fingerprints). The embeddings already summarize the text the
labels were written from. The other fields describe the trial, and year points
the wrong way because the labels are early and the test is late.

Missing labels have to come from an unsupervised step. $k$-means and DBSCAN
were useful checks, and one DBSCAN-plus-logistic setup reached about 0.910
validation F1 under an earlier protocol. We did not submit it. Clustering
mostly separates Recruitment from everything else. The piece we kept is a
30-component diagonal Gaussian mixture on 30 principal components, about 77\%
of the variance~\cite{dempster1977em,pedregosa2011scikit}. Components start from
$k$-means++, three restarts, \texttt{random\_state}${}=0$. Expectation--maximization
stops when the average log-likelihood improves by less than $10^{-3}$, or after
100 iterations. The kept run converged in 29. Each component becomes a reason
profile by averaging the real labels assigned to it, plus one pseudo-count of
the overall rates. In the teaching step this profile has a 5\% vote. The
classifier has the other 95\%.

The classifier is one logistic regression over the five reasons ($C = 0.1$),
not five separate classifiers. A two-reason trial is entered once per reason,
with its weight split, so it is not counted twice. Class weights come from the
real labels only and then stay fixed. A generated row counts as a quarter of a
real row. We average three scores with fixed weights: logistic regression 50\%,
cosine 10-nearest neighbors 30\%, and shrinkage linear discriminant analysis
20\%. Neighbors vote with squared cosine similarity, and a neighbor that has
two reasons votes for both. Only the mixture uses the 30 components. The
classifiers see all 150 standardized embeddings.

Unlabeled trials are added in two rounds~\cite{yarowsky1995unsupervised,triguero2015selflabeled}.
A trial is accepted when its best reason scores at least 0.8. Each round
starts again from the real labels. On the final fit the rounds accept 685 and
then 757 trials, with total weight about 103 and then 120, against 235 for the
real labels. We always output the top reason, plus any other reason at 0.3 or
above, because an empty prediction is always wrong. Test rows are used only for
the final predictions. The 0.891 validation score belongs to a fit on 157
labels. The submitted file is retrained on all 235, so that validation number
is not an independent score of the submitted file. An earlier 0.89423
submission had fit its scaler and mixture with the test rows included. We
stopped doing that.

\section{Part 2: trial outcomes}
\label{sec:p2}

\subsection{Data and the score}

Each row is one trial in one phase. A combined-phase trial appears twice and
those rows stay together in every split. Training has 8,037 rows through 2013,
validation has 1,013 rows from the same years, and the test file has 1,500
later rows, 500 per phase, from 2013 to 2022. Kaggle averages the three phase
log losses. We select on that same average. F1 at 0.5, ROC-AUC, and
precision-recall AUC are diagnostics. Part~1 reasons are joined only after
prediction, for the error table. They are not model inputs.

\subsection{What the model does}

The baseline is logistic regression on eligibility text: word and bigram
counts, at most 35,000 terms, minimum document frequency 3, sublinear counts,
three phase indicators, and $C = 2$. Separate phase-only fits overfit Phase~I.
Drug fingerprints did not reliably help, and a matched drug is less common in
test (39\%) than in training (46\%). Naive Bayes ranked trials but was far too
confident for log loss.

The model we built on is a stack of eight classifiers~\cite{wolpert1992stacked}:
logistic models on eligibility text, diseases, drug names, codes, drug
annotations, and metadata, plus naive Bayes and a 50-neighbor model on the
eligibility text. A second logistic regression learns the weights from
out-of-fold predictions, five folds, with a trial never split across folds.
Recent trials get more weight, with a three-year exponential time constant, so
a trial loses about half its weight every two years. That stack, refit on all
9,050 labeled rows, scored 0.58688. The same stack without recency weights
scored 0.58810. Our first ensembles had used only the 8,037 training rows, so
part of the earlier gap was the extra data.

The neighbor piece inside the stack ignored time and pulled estimates toward a
50\% success rate. File~10 does not replace the stack. It averages it, half and
half, with a second stack that decays faster (two-year time constant) and that
also sees eight neighborhood estimates: several sizes, sharper similarity, and
similarity that includes disease, drug, and ICD text. The weights are fixed.
This improved every time-ordered window and scored 0.58609, even though the
random validation log loss got slightly worse.

File~10 is weaker when the years are mixed. Our earlier ensemble is weaker on
later years. That ensemble is half three logistic models, a quarter a depth-6
tree, and a quarter a depth-8 tree. We first added 20\% of it (Kaggle 0.58474).
Raising that share to 30\% gained only another 0.00009, so we stopped. The
submitted probability is
\[
0.70\,P_{10} + 0.30\,P_{\text{early}},
\quad
P_{10} = 0.50\,P_{\text{recency}} + 0.50\,P_{\text{fast neighbors}}.
\]
Supervised pieces are refit on all 9,050 labeled rows after the recipe is
frozen. Test labels are never used. In the inherited stack, word lists may see
unlabeled test text. That changes the features, not the answers. A later blend
with an extra tree and a phase trend looked better locally and scored 0.58492,
so we did not replace file~13. One upload slot was left unused.

\section{Results}
\label{sec:results}

Table~\ref{tab:p1} is why Part~1 submitted the final blend rather than the
linear-discriminant ensemble. The two are tied under cross-validation. The
ensemble wins the 78-trial validation score. The final blend wins the
time-ordered check and the reported Kaggle score. On those 78 trials its
per-reason F1 is 0.941 Safety (9 trials), 0.667 Funding (4), 0.759 Efficacy
(14), 0.990 Recruitment (50), and 0.471 Other (10): sample-wise 0.891, macro
0.765.

\begin{table}[t]
\centering\small
\setlength{\tabcolsep}{3pt}
\begin{tabular*}{\columnwidth}{@{\extracolsep{\fill}}lccccc@{}}
\toprule
Model & Train & Pooled & Val. & Time & Kaggle \\
\midrule
Final blend & \textbf{.9321} & .9214 & .891 & \textbf{.884} & \textbf{.89636} \\
LDA ensemble & .9318 & \textbf{.9216} & \textbf{.906} & .878 & .89401 \\
Earlier LDA blend & .9281 & .9154 & .872 & .874 & --- \\
First submission & .9196 & .9098 & .891 & .861 & .89423 \\
\bottomrule
\end{tabular*}
\caption{Part~1 on shared grouped checks. Train uses 157 labels, pooled uses
235, validation is the original 78 trials, and time pools three chronological
windows. Higher is better.}
\label{tab:p1}
\end{table}

Table~\ref{tab:p2} is the Part~2 path. Lower log loss is better. File~13 is the
submission. File~15 wins both local columns and loses on Kaggle, which is why
the local ranking was not enough.

\begin{table}[t]
\centering\small
\setlength{\tabcolsep}{3pt}
\begin{tabular*}{\columnwidth}{@{\extracolsep{\fill}}lccc@{}}
\toprule
Model & Time & Grouped & Kaggle \\
\midrule
Phase-specific ensemble & --- & --- & 0.59562 \\
Text-heavy ensemble & --- & --- & 0.58993 \\
Stack, no recency & 0.5421 & 0.5797 & 0.58810 \\
Recency stack & 0.5374 & 0.5874 & 0.58688 \\
File 10, recent neighbors & 0.5359 & 0.5869 & 0.58609 \\
File 13, 30\% early ensemble & 0.5334 & 0.5790 & \textbf{0.58465} \\
File 15, tree and trend & \textbf{0.5327} & \textbf{0.5783} & 0.58492 \\
\bottomrule
\end{tabular*}
\caption{Part~2 log loss. Time averages the forward windows. Grouped is
trial-grouped cross-validation. The first two submissions used an earlier
local protocol, so those cells are blank.}
\label{tab:p2}
\end{table}

Table~\ref{tab:phase} is the supplied validation split, with preprocessing fit
only on that split's training rows. The recency stack improves Phase~III log
loss from 0.529 to 0.516 and Phase~I F1 from 0.765 to 0.793. It does not
improve Phase~II F1. File~10 is slightly worse than the recency stack on this
random split and better on the later-year windows in Table~\ref{tab:p2}. That
is why a higher validation F1 was not the submission rule.

\begin{table}[t]
\centering\small
\setlength{\tabcolsep}{2.5pt}
\begin{tabular}{llccc}
\toprule
Phase & Metric & Baseline & Recency & File 10 \\
\midrule
\multirow{4}{*}{I}
 & F1       & 0.765 & \textbf{0.793} & 0.762 \\
 & ROC-AUC  & 0.783 & \textbf{0.794} & 0.787 \\
 & PR-AUC   & \textbf{0.835} & 0.830 & 0.824 \\
 & log loss & 0.556 & 0.556 & 0.560 \\
\midrule
\multirow{4}{*}{II}
 & F1       & \textbf{0.657} & 0.647 & 0.638 \\
 & ROC-AUC  & 0.740 & 0.753 & \textbf{0.754} \\
 & PR-AUC   & 0.717 & \textbf{0.732} & 0.730 \\
 & log loss & 0.600 & 0.593 & 0.593 \\
\midrule
\multirow{4}{*}{III}
 & F1       & \textbf{0.833} & 0.825 & 0.829 \\
 & ROC-AUC  & 0.774 & \textbf{0.790} & 0.786 \\
 & PR-AUC   & 0.862 & \textbf{0.882} & 0.877 \\
 & log loss & 0.529 & \textbf{0.516} & 0.521 \\
\bottomrule
\end{tabular}
\caption{Part~2 validation by phase. F1 uses a cutoff of 0.5. The other rows
use the probability of success. Phase sizes are 149, 501, and 363. Best value
per row in bold. Tied log losses are left unbolded.}
\label{tab:phase}
\end{table}

Table~\ref{tab:err} counts failed validation trials that the outcome model
called successful. The ensemble is the early text-and-tree model that later
became 30\% of file~13. The readable groups are Phase~II Recruitment and the
three no-reason groups. Phase~III failures are called successful much more
often than Phase~II failures. Most reason tags are Part~1 predictions. The
later neighbor model still called all 9 Phase~III efficacy-tagged failures
successful.

\begin{table}[t]
\centering\small
\setlength{\tabcolsep}{3pt}
\begin{tabular*}{\columnwidth}{@{\extracolsep{\fill}}llccc@{}}
\toprule
Phase & Reason & $n$ & Base & Ens. \\
\midrule
I & Safety & 9 & .56 & .44 \\
I & Funding & 4 & .50 & .50 \\
I & Efficacy & 2 & .50 & .50 \\
I & Recruitment & 22 & .36 & .27 \\
I & Other & 5 & .80 & .80 \\
I & No reason & 28 & .43 & .39 \\
II & Safety & 15 & .40 & .40 \\
II & Funding & 4 & .00 & .25 \\
II & Efficacy & 10 & .30 & .30 \\
II & Recruitment & 87 & .20 & .16 \\
II & Other & 18 & .33 & .28 \\
II & No reason & 144 & .35 & .34 \\
III & Safety & 6 & .50 & .67 \\
III & Funding & 0 & --- & --- \\
III & Efficacy & 9 & .89 & .67 \\
III & Recruitment & 21 & .52 & .57 \\
III & Other & 8 & .63 & .75 \\
III & No reason & 82 & .56 & .55 \\
\bottomrule
\end{tabular*}
\caption{Share of failed validation trials predicted successful. ``No reason''
means the trial is absent from Part~1. A trial with several reasons is listed
under each.}
\label{tab:err}
\end{table}

\section{Why these choices}
\label{sec:why}

The Part~1 mixture is a poor labeler and a useful tie-breaker. On its own it
was the weakest model we tried, because extra components split Recruitment
rather than isolating Safety or Funding. A supervised model still reaches 0.941
F1 on Safety from the same embeddings, so the features are not the problem.
Cutting the mixture's vote to 5\% and the generated-row weight to one quarter
lets a confident guess help without outvoting the 235 real labels.

Extra Part~1 features fail for a mechanical reason. The embeddings are a
compressed copy of the text the labels came from. Criteria, diseases, and drugs
describe what the trial was. Adding them means dozens or thousands of new
inputs on 157 real labels, and 44\% of trials have no matched drug. Start year
looks predictive only because the labeled trials were chosen to be early.

We did not follow Part~1's higher validation score. Four Funding trials and ten
Other trials can move the macro average by several points. The time-ordered
windows look more like the test years, and they favor the submitted blend, as
does the reported Kaggle score.

In Part~2, log loss and F1 disagree on purpose. A model that moves an uncertain
trial across 0.5 can lose F1 and still assign the better probability. The test
file is a later, more successful cohort than the random validation split, so we
trusted time-ordered checks. Recency weights and recent neighbors are two ways
of saying that an old similar trial is a weaker example. The 30\% older
ensemble is there because it misses different trials, and the second 10\% of
that share was only worth 0.00009 on Kaggle. After file~15 lost on the
leaderboard, a better historical average was no longer a reason to upload.

The reason table says where outcome errors live, not how to tune a cutoff.
Failed Phase~III trials often ran to the end and simply missed the endpoint, so
the eligibility text looks like a success. Recruitment and funding failures are
more administrative, and the ensemble catches a few more of them. Safety,
Funding, and Other are small in every phase. Only Phase~II Recruitment and the
no-reason groups are large, and the no-reason group is a lookup rather than a
prediction.

{
    \small
    \bibliographystyle{midterm/report/ieeenat_fullname}
    \bibliography{midterm/report/main}
}

\end{document}
